\subsection{分次素理想}

设 \(A = \bigoplus_{i \geqslant 0} A_i\) 是一个正次数的分次交换环，\(m\) 是分次理想 \(\bigoplus_{i \geqslant 1} A_i\) ；我们称 A 的两个分次理想 \(a = \bigoplus_{i \geqslant 0} a_i\) 、\(b = \bigoplus_{i \geqslant 0} b_i\) 是等价的，如果存在一个整数 \(n_0\) ，使得 \(a_n = b_n\) 对 \(n \geqslant n_0\) 成立（显然这是一个等价关系）。一个分次理想称为本性的，如果它不等价于 \(m\) 。

命题 4. 设 \(\mathfrak{p} = \bigoplus_{i\geqslant 0}\mathfrak{p}_i\) 是 A 的一个分次理想；为使 \(\mathfrak{p}\) 是素的，必须且只需：若 \(x\in \mathbf{A}_m,y\in \mathbf{A}_n\) 满足 \(x\notin \mathfrak{p}\) 且 \(y\notin \mathfrak{p}\) ，则 \(xy\notin \mathfrak{p}\) 。

该条件显然是必要的。反之，若它成立，则在分次环 \(A/p = \bigoplus_{i \geq 0} A_i/p_i\) 中，两个 \(\neq 0\) 的齐次元之积 \(\neq 0\) ，因而 A/p 是一个整环（《代数》第 II 章 §11 第 4 节命题 7）。

命题 5. 设 \(\mathfrak{a} = \bigoplus_{i\geqslant 0}\mathfrak{a}_i\) 是 A 的一个分次理想，\(n_0\) 是一个 \(>0\) 的整数。为使存在一个分次素理想 \(\mathfrak{p} = \bigoplus_{i\geqslant 0}\mathfrak{p}_i\) ，使得 \(\mathfrak{p}_n = \mathfrak{a}_n\) 对 \(n\geqslant n_0\) 成立，必须且只需：对齐次元素对 \(x,y\) （次数 \(\geqslant n_0\) ），关系 \(xy\in \mathfrak{a}\) 蕴含“ \(x\in \mathfrak{a}\) 或 \(y\in \mathfrak{a}"\) 。若存在 \(n\geqslant n_0\) 使 \(\mathfrak{a}_n\neq \mathbf{A}_n\) ，则满足上述条件的素理想是唯一的。

陈述中的条件显然是必要的。若 \(a_{n} = A_{n}\) 对所有 \(n \geqslant n_{0}\) 成立，则显然每个包含 m 的素理想都是该问题的一个解；因而可能有多个素理想解决该问题；不过这些理想中的任意两个显然等价。于是设存在一个齐次元 \(a \in A_{d}\) （其中 \(d \geqslant n_{0}\) ），它不属于 \(a_{d}\) 。设 p 是由 \(x \in A\) （满足 \(ax \in a\) ）组成的集合。显然 p 是 A 的一个理想；由于 ax 的齐次分量是 a 与 x 的诸齐次分量的乘积，且 a 是分次理想，故 p 是分次理想；此外，\(1 \notin p\) ，因而 \(p \neq A\) 。为证明 p 是素的，只需证明：若 \(x \in A_{m}\) 、\(y \in A_{n}\) 满足 \(x \notin p\) 与 \(y \notin p\) ，则 \(xy \notin p\) （命题 4）。这时 \(ax \notin a_{m+d}\) ，\(ay \notin a_{n+d}\) ，从而由假设

\[
a ^ {2d} x y \notin \mathfrak {a} _ {m + n + 2 d};
\]

我们断言 \(axy \notin \mathfrak{a}_{m+n+d}\) ，因为 \(xy \notin \mathfrak{p}\) 。最后，若 \(n \geqslant n_0\) 且 \(x \in \mathbf{A}_n\) ，则条件 \(x \in \mathfrak{a}_n\) 与 \(ax \in \mathfrak{a}_{n+d}\) 依假设等价，因而 \(\mathfrak{p} \cap \mathbf{A}_n = \mathfrak{a}_n\) ，这就完成了对该问题的解——分次素理想 \(\mathfrak{p}\) 的存在性证明。若 \(\mathfrak{p}'\) 是 \(\mathbf{A}\) 的另一个分次素理想，满足 \(\mathfrak{p}' \cap \mathbf{A}_n = \mathfrak{a}_n\) 对 \(n \geqslant n_0\) 成立，则 \(a \notin \mathfrak{p}'\) ，且 \(ax \in \mathfrak{p}'\) 对所有 \(x \in \mathfrak{p}\) 成立，从而 \(\mathfrak{p} \subset \mathfrak{p}'\) ，因为 \(\mathfrak{p}'\) 是素的。另一方面，若 \(x\) 是一个次数为 \(n \geqslant 0\) 的齐次元，属于 \(\mathfrak{p}'\) ，则 \(ax\) 是次数为 \(n + d \geqslant n_0\) 的齐次元，因此属于 \(\mathfrak{p}' \cap \mathrm{A}_{n+d} = \mathfrak{a}_{n+d}\) ，从而依定义 \(x \in \mathfrak{p}\) ，这表明 \(\mathfrak{p}' \subset \mathfrak{p}\) ，最终 \(\mathfrak{p}' = \mathfrak{p}\) 。

命题 6. 设 \(d\) 是一个 \(\geqslant 1\) 的整数。

\begin{enumerate}
\def\labelenumi{(\roman{enumi})}
\item
  对每个本性分次理想 \(\mathfrak{p}\) ，\(\mathbf{A},\mathfrak{p}\cap \mathbf{A}^{(d)}\) 是 \(\mathbf{A}^{(d)}\) 的一个本性分次素理想。
\item
  反之，对每个本性分次素理想 \(\mathfrak{p}'\) （属于 \(\mathbf{A}^{(d)}\) ），存在唯一的（必然是本性的）分次素理想 \(\mathfrak{p}\) （属于 \(\mathbf{A}\) ），使得 \(\mathfrak{p} \cap \mathbf{A}^{(d)} = \mathfrak{p}'\) 。
\item
  若 \(a \in \mathbf{A}_k\) 不属于 \(\mathfrak{p}_k\) ，则 \(a^{kd}\) 不属于 \(\mathfrak{p}_{kd}\) ，因而 \(\mathfrak{p} \cap \mathbf{A}^{(d)}\) 是本性的。
\item
  对所有 \(n \geqslant 0\) ，集合 \(\mathfrak{p} \cap \mathbf{A}_n\) 必须等于集合 \(\mathfrak{a}_n\) ，后者由 \(x \in \mathbf{A}_n\) （满足 \(x^d \in \mathfrak{p}'\) ）组成。我们来证明 \(\mathfrak{a} = \bigoplus_{n \geqslant 0} \mathfrak{a}_n\) 是一个分次素理想；由于 \(\mathfrak{a}_n = \mathfrak{p}_n'\) 当 \(n\) 是 \(d\) 的倍数时成立（因为 \(\mathfrak{p}'\) 是素的），这就将证明 \(\mathfrak{p}\) 的唯一性。现在，若 \(x \in \mathfrak{a}_n, y \in \mathfrak{a}_n, (x - y)^{2d}\) 是一些项之和，其中每一项都是 \(x^d\) 或 \(y^d\) 与一个次数为 \(nd\) 的齐次元的乘积，因而 \((x - y)^{2d} \in \mathfrak{p}'\) ，且由于 \(\mathfrak{p}'\) 是素的，\((x - y)^d \in \mathfrak{p}'\) ，因此 \(\mathfrak{a}_n\) 是 \(\mathbf{A}\) 的一个子群。由于 \(\mathfrak{p}'\) 是 \(\mathbf{A}^{(d)}\) 的理想，\(\mathfrak{a}\) 是 \(\mathbf{A}\) 的分次理想；最后，关系 \((xy)^d \in \mathfrak{p}'\) 蕴含 \(x^d \in \mathfrak{p}'\) 或 \(y^d \in \mathfrak{p}'\) ，依命题 4，这就完成了证明。
\end{enumerate}

设 A 是一个正次数的分次交换环，p 是 A 的一个本性分次素理想。A 中不属于 p 的齐次元组成的集合 S 是乘性的，因而分式环 \(S^{-1}A\) 典范地是分次的（第 II 章 §2 第 9 节）（注意在这个分次中一般会存在次数为负的 \(\neq0\) 的齐次元）。我们将用 \(\mathrm{A}_{(\mathrm{p})}\) 表示 \(S^{-1}A\) 的由次数为 0 的齐次元组成的子环，换言之，即形如 x/s 的分式组成的集合，其中 x 与 s 是 A 中次数相同的齐次元，且 \(s\notin p\) 。类似地，对每个分次 A-模 M，\(S^{-1}M\) 典范地是分次的（出处同上），我们将用 \(\mathrm{M}_{(\mathrm{p})}\) 表示次数为 0 的齐次元组成的子群，它显然是一个 \(\mathrm{A}_{(\mathrm{p})}\) -模。

命题 7. 设 \(\mathfrak{p}\) 是 \(\mathbf{A}\) 的一个分次素理想，\(d\) 是一个 \(\geqslant 1\) 的整数，\(\mathfrak{p}'\) 是分次素理想 \(\mathfrak{p} \cap \mathrm{A}^{(d)}\) （属于 \(\mathrm{A}^{(d)}\) ）；对每个分次 A-模 \(\mathbf{M}\) ，同态 \((\mathbf{M}^{(d)})_{(\mathfrak{p}')} \to \mathbf{M}_{(\mathfrak{p})}\) （由典范单射 \(\mathbf{M}^{(d)} \to \mathbf{M}\) 导出）是双射。

若 S 是 A 中不属于 p 的齐次元组成的集合，且 \(\mathrm{S}^{\prime}=\mathrm{S}\cap\mathrm{A}^{(d)}\) ，则典范同态 \(\phi:\mathrm{S}^{\prime-1}\mathrm{M}^{(d)}\to\mathrm{S}^{-1}\mathrm{M}\) 是一个次数为 0 的齐次同态，并且是单射，因为若 \(x\in M_{nd}\) 对 \(s\in A_{m}\) 满足 sx=0，这里 \(s\notin p\) ，则也有 \(s^{d}x=0\) ，且 \(s^{d}\in A_{md}\) ，\(s^{d}\notin p^{\prime}\) 。剩下要证 \(\phi\) 把 \((\mathrm{M}^{(d)})_{(\mathfrak{p}^{\prime})}\) 映到整个 \(\mathrm{M}_{(\mathfrak{p})}\) ；但若 \(x\in M_{n}\) ，\(s\in A_{n}\) 且 \(s\notin p\) ，则也有 \(x/s=(xs^{d-1})/s^{d}\) ，其中 \(xs^{d-1}\in A_{nd}\) ，\(s^{d}\in A_{nd}\) ，且 \(s^{d}\notin p^{\prime}\) ，由此即得我们的断言。

命题 8. 设 \(\mathfrak{m} = \bigoplus_{i\geqslant 1}\mathbf{A}_i\) ；设 \((\mathfrak{p}^{(k)})_{1\leqslant k\leqslant n}\) 是 A 的分次素理想的一个有限族，a 是 A 的一个分次理想，使得对所有 k 有 \(a \cap m \notin p^{(k)}\) ；则存在一个齐次元 \(z \in a \cap m\) ，它不属于任何 \(p^{(k)}\) 。

我们对 \(n\) 作归纳论证，命题对 \(n = 1\) 是平凡的。若存在一个指标 \(j\) ，使得 \(\mathfrak{a} \cap \mathfrak{m} \cap \mathfrak{p}^{(j)}\) 包含于诸 \(\mathfrak{p}^{(k)}\) （指标 \(k \neq j\) ）之一，则由归纳假设，存在一个齐次元 \(z' \in \mathfrak{a} \cap \mathfrak{m}\) ，它不属于诸 \(\mathfrak{p}^{(k)}\) （指标 \(k \neq j\) ）中的任何一个，因而不属于 \(\mathfrak{p}^{(j)}\) ，这个元素就解决了该问题。于是设对每个指标 \(j\) ，\(\mathfrak{a} \cap \mathfrak{m} \cap \mathfrak{p}^{(j)}\) 都不包含于诸 \(\mathfrak{p}^{(k)}\) （指标 \(k \neq j\) ）中的任何一个；因此归纳假设蕴含存在一个齐次元 \(y_j \in \mathfrak{a} \cap \mathfrak{m} \cap \mathfrak{p}^{(j)}\) ，它不属于诸 \(\mathfrak{p}^{(k)}\) （指标 \(k \neq j\) ）中的任何一个；由于诸 \(y_j\) 的次数都 \(\geqslant 1\) ，我们可以假设（因为诸 \(\mathfrak{p}^{(k)}\) 都是素的，用适当的幂代替它们）\(y_1\) 与 \(\prod_{j=2}^{n} y_j\) 的次数相同。于是 \(z = y_1 + \prod_{j=2}^{n} y_j\) 是次数 \(\geqslant 1\) 的齐次元，并且与第 II 章 § 1 第 1 节命题 2 中相同的论证表明 \(z\) 解决了该问题。

